\section{Generador de órdenes médicas}

El generador produce órdenes médicas periódicas cuyo caudal e intervalo
de tiempo se muestrean de las distribuciones definidas en el documento
de modelado estadístico, parametrizadas por el grado de criticidad $k$.
El parámetro $k$ es una constante del modelo que representa el grado
de criticidad asignado al paciente y no varía durante la simulación.

\subsection{Definición formal}

Sea $k_0 \in [0,1]$ el grado de criticidad del paciente, fijado como
parámetro del modelo.

\[
  M_{\mathrm{gen}} = \langle X,\; Y,\; S,\; \ta,\; \dint,\; \dext,\; \lambda \rangle
\]

\begin{align*}
  X &= \varnothing \\[4pt]
  Y &= \bigl\{\,(c,\,t,\;0)\;\big|\;
        c\in\{0,50,100,150,200\},\;
        t\in\{2,4,6,8,12\}\bigr\} \\[4pt]
  S &= \bigl\{\,(c,\,t)\;\big|\;
        c\in\{0,50,100,150,200\},\;
        t\in\{2,4,6,8,12\}\bigr\}
\end{align*}

El parámetro $k_0$ no forma parte del estado sino que es una constante
del modelo que parametriza las funciones $\mathrm{Caudal}$ y $\mathrm{Tiempo}$.

\paragraph{Transición interna.}
Genera una nueva orden a partir del grado de criticidad del paciente:
\[
  \dint(c,\,t)
  \;=\;
  \bigl(\,\mathrm{Caudal}(k_0),\;\mathrm{Tiempo}(k_0)\,\bigr).
\]

\paragraph{Transición externa.}
El modelo no recibe eventos externos:
\[
  \dext = \varnothing.
\]

\paragraph{Función de salida.}
Emite el caudal y el tiempo de la orden vigente:
\[
  \lambda(c,\,t) \;=\; (c,\,t,\;0).
\]

\paragraph{Avance temporal.}
\[
  \ta(c,\,t) \;=\; t.
\]

\paragraph{Estado inicial.}
\[
  s_0 \;=\; \bigl(\,\mathrm{Caudal}(k_0),\;\mathrm{Tiempo}(k_0)\,\bigr),
\]
donde $k_0$ es el grado de criticidad del paciente, fijado como parámetro
al instanciar el modelo.

\subsection{Diagrama de estados}

\begin{figure}[htbp]
\centering
\begin{tikzpicture}
\node[draw,circle] (gen) {$\left(c,t\right)$};
\node[draw,rectangle] (ctrl) at (5,0) {Controlador};

\draw[-{Stealth}] (gen) -- (ctrl)
node[midway,above] {$\lambda=(c,t,0)$};

\draw[-{Stealth}]
(gen) edge[loop left]
node {$\dint(c,t) = (\mathrm{Caudal}(k_0),\mathrm{Tiempo}(k_0))$}
(gen);
\end{tikzpicture}
\caption{Diagrama de estados del Generador de Órdenes Médicas.}
\label{fig:diag_generador_med}
\end{figure}